% ===== PRÉAMBULE CANONIQUE — à coller VERBATIM en tête de CHAQUE .tex =====
\documentclass[11pt,a4paper]{article}
\usepackage[utf8]{inputenc}\usepackage[T1]{fontenc}\usepackage{lmodern}
\usepackage[french]{babel}
\usepackage{amsmath,amssymb,mathtools}\usepackage{icomma}
\usepackage[version=4]{mhchem}          % \ce{...} : équations & formules
\usepackage{siunitx}\sisetup{locale=FR} % \qty{}{}, \si{}, \ang{}
\usepackage{chemfig}                    % \chemfig{...} : structures développées
\usepackage{xcolor}\usepackage{colortbl}
\usepackage{tikz}\usepackage{pgfplots}\pgfplotsset{compat=1.18}
\usepackage[most]{tcolorbox}\usepackage{enumitem}\usepackage{array,booktabs}
\usepackage[a4paper,margin=2.2cm]{geometry}
\usetikzlibrary{arrows.meta,calc,angles,quotes,positioning,patterns,backgrounds,shapes.geometric}
\definecolor{primary}{RGB}{222,49,99}\definecolor{primarydark}{RGB}{150,22,55}
\definecolor{defcol}{RGB}{0,114,178}\definecolor{methcol}{RGB}{52,92,148}
\definecolor{excol}{RGB}{0,135,90}\definecolor{attncol}{RGB}{200,40,55}
\tcbset{boxbase/.style={enhanced,breakable,boxrule=0.9pt,arc=2.5pt,
  left=3mm,right=3mm,top=2.3mm,bottom=2.3mm,fonttitle=\bfseries,coltitle=white}}
% --- boîtes TUTORIEL (noms EXACTS, ne pas renommer) ---
\newtcolorbox{cle}[1][]{boxbase,colback=defcol!6,colframe=defcol,
  title={\textsc{À retenir}\ifx\relax#1\relax\relax\else\textmd{ : \itshape #1}\fi}}
\newtcolorbox{methode}[1][]{boxbase,colback=methcol!7,colframe=methcol,sharp corners,
  title={\textsc{Méthode}\ifx\relax#1\relax\relax\else\textmd{ --- \itshape #1}\fi}}
\newtcolorbox{exemplebox}[1][]{boxbase,colback=excol!5,colframe=excol,
  title={\textsc{Exemple}\ifx\relax#1\relax\relax\else\textmd{ : \itshape #1}\fi}}
\newtcolorbox{piege}[1][]{boxbase,colback=attncol!6,colframe=attncol,
  title={\textsc{Piège}\ifx\relax#1\relax\relax\else\textmd{ : \itshape #1}\fi}}
\newtcolorbox{remarque}[1][]{boxbase,colback=black!4,colframe=black!45,
  title={\textsc{Remarque}\ifx\relax#1\relax\relax\else\textmd{ : \itshape #1}\fi}}
% --- boîtes EXERCICE / CORRIGÉ (noms EXACTS, auto-comptées) ---
\newtcolorbox[auto counter]{exo}[1][]{boxbase,colback=primary!4,colframe=primary,
  title={\textsc{Exercice}~\thetcbcounter\ifx\relax#1\relax\relax\else\textmd{ --- \itshape #1}\fi}}
\newtcolorbox[auto counter]{corrige}[1][]{boxbase,colback=excol!4,colframe=excol,
  title={\textsc{Corrigé}~\thetcbcounter\ifx\relax#1\relax\relax\else\textmd{ --- \itshape #1}\fi}}
\newcommand{\R}{\mathbb{R}}\newcommand{\N}{\mathbb{N}}\newcommand{\Z}{\mathbb{Z}}\newcommand{\ds}{\displaystyle}
% (un \usepackage{fancyhdr}+en-tête est possible ; il est ignoré côté web)
% =========================================================================
\begin{document}

% THEME_TITLE: Dilution, fractions et molalité
\noindent
Ce thème regroupe les opérations et grandeurs « relatives » des solutions : la \textbf{dilution}
(ajouter du solvant) et les descripteurs d'un \textbf{mélange} (fractions, molalité).

\begin{cle}[dilution : la quantité de soluté se conserve]
Diluer, c'est ajouter du solvant \emph{sans toucher au soluté} : $n$ est conservée. En notant $1$ la
solution \emph{mère} (avant) et $2$ la \emph{fille} (après) :
\[ \boxed{\,C_1 V_1 = C_2 V_2\,}. \]
Comme on ajoute de l'eau, $V_2 > V_1$ donc $C_2 < C_1$ : \textbf{une dilution ne peut que diminuer la
concentration}. Le \textbf{facteur de dilution} est $F = \dfrac{V_2}{V_1} = \dfrac{C_1}{C_2}$
(« diluer $5$ fois » $\Rightarrow F=5$).
\end{cle}

\begin{center}
\begin{tikzpicture}[scale=1,bk/.style={draw,line width=0.9pt,rounded corners=1pt}]
  \fill[primary!30] (0,0) rectangle (1.1,1.0); \draw[bk] (0,0) rectangle (1.1,1.6);
  \node[font=\scriptsize] at (0.55,-0.35) {mère $C_1,V_1$};
  \draw[-{Stealth[length=2mm]},line width=0.9pt,defcol] (1.5,0.8) -- node[above,font=\scriptsize]{$+\,\ce{H2O}$} (3.0,0.8);
  \fill[primary!12] (3.4,0) rectangle (5.0,1.3); \draw[bk] (3.4,0) rectangle (5.0,1.9);
  \node[font=\scriptsize] at (4.2,-0.35) {fille $C_2<C_1,\ V_2>V_1$};
  \node[draw=primarydark,rounded corners=2pt,fill=primarydark!6,font=\footnotesize] at (8.0,0.8) {$C_1V_1=C_2V_2$};
\end{tikzpicture}
\end{center}

\begin{cle}[fraction molaire et fraction massique]
Pour un mélange de constituants \ce{A}, \ce{B}\ldots :
\[ \chi_{\ce{A}} = \dfrac{n_{\ce{A}}}{n_{\ce{A}}+n_{\ce{B}}+\cdots}\qquad
   w_{\ce{B}} = \dfrac{m_{\ce{B}}}{m_{\ce{A}}+m_{\ce{B}}+\cdots}. \]
Toutes deux sont \textbf{sans unité}, entre $0$ et $1$, et $\sum\chi = 1$ (de même $\sum w = 1$).
Le \textbf{pourcentage massique} de \ce{B} vaut $w_{\ce{B}}\times 100$.
\end{cle}

\begin{cle}[molalité]
La \textbf{molalité} $m$ est le nombre de moles de soluté par \textbf{kilogramme de solvant} :
\[ \boxed{\,m = \dfrac{n(\text{soluté})}{m_{\text{solvant}}\ (\text{en \si{\kilo\gram}})}\,},\qquad
   \text{unité : }\si{\mole\per\kilo\gram}. \]
\end{cle}

\begin{piege}[erreurs classiques]
\begin{itemize}[leftmargin=1.4em,topsep=1pt,itemsep=1pt]
  \item On ne peut \textbf{jamais} concentrer une solution en la diluant ($C_2 < C_1$).
  \item La \textbf{molalité} se rapporte à la masse de \emph{solvant} en \textbf{kg} (pas en g, pas la
        solution entière) ; à ne pas confondre avec la molarité (\si{\mole\per\litre}).
\end{itemize}
\end{piege}

\end{document}
